% Checked algebraic deduction; whole-proof archive audit in cup_ext_archive_audit.md.
% This lemma alone does not assert a kernel theorem without its module hypotheses.
\subsection{Finite-level cups detect the actual companion poles}

The following comparison concerns the actual class-two extension of
the Galois group. It supplies the naturality needed to use a cup-product
theorem; an equality of characteristic ideals would not supply it.

\begin{lemma}[K17-CUP-EXT]\label{lem:cup-pole-bridge}
Let $p$ be odd, let $P=H\rtimes\Gamma$, where $H$ is free pro-$p$
and $\Gamma=\langle\gamma\rangle\simeq\mathbb Z_p$, and let an
involution $c$ preserve $H$ and fix $\gamma$. Put
\[
 \Lambda=\mathbb Z_p[[\Gamma]],\qquad R=\Lambda[1/p],\qquad
 H^{\rm ab}=U\oplus T,
\]
where $c=-1$ on $U$ and $c=1$ on $T$. Suppose that $U$ is finitely
generated over $\Lambda$, that $U[1/p]$ is free of finite rank over
$R$, and that $T$ is finite free over $\mathbb Z_p$. Suppose also
that each even character component $T_e[1/p]$ under the commuting
prime-to-$p$ action is cyclic over $R$. Write $q_e(z)$ for its monic
characteristic and minimal polynomial for $\gamma$.

Let the coefficient character satisfy $\chi(\gamma)=u\in1+p\mathbb Z_p$
and $\chi(c)=-1$, with $H$ acting trivially on its twists. Assume that,
for every $n\geq1$, the even part of
\[
 H^2(P_n,\mathbb Z/p^n(2)),\qquad
 P_n=H\rtimes\Gamma^{p^n},
\]
is generated by cups of $c$-even classes in
$H^1(P_n,\mathbb Z/p^n(1))$. A cofinal diagonal of the two level
indices suffices.

Then the actual odd--odd companion numerators satisfy the
pole-nonvanishing condition: for every root $\alpha$ of every
nonzero $q_e$, at least one coefficient has nonzero restriction to
\[
 D_\alpha=(1+T_1)(1+T_2)-\alpha=0.
                                                               \tag{CE.1}
\]
Here coefficients are taken in the bounded-denominator power-series
coordinates of an $R$-basis of $U[1/p]$. Repeated roots are allowed.
For the subsequent pole-separation application the arithmetic module
input supplies rank one in every odd character.
\end{lemma}

\begin{proof}
Since $2$ is a unit, the class-two quotient of $H$ has its integral
Baker--Campbell--Hausdorff Lie structure. Quotient its central
commutator module by all brackets having a $T$ input. The remaining
central module is
\[
 N=\widehat{\bigwedge}^{\,2}_{\mathbb Z_p}U.
\]
Complex conjugation is $1$ on $N$. Thus the minus part of the
resulting Lie module maps isomorphically to $U$, providing a
canonical $\Gamma$-equivariant odd section. Its plus part is the
continuous module extension
\[
 0\longrightarrow N\longrightarrow E\longrightarrow T
 \longrightarrow0.                                          \tag{CE.2}
\]
Choose a rational cyclic generator $v$ of $T_e[1/p]$ and a lift
$y\in E[1/p]$ in the same character component, using the
prime-to-$p$ character projector. The element
\[
 q_e(\gamma)y=C_e\in N[1/p]                                 \tag{CE.3}
\]
is precisely the odd--odd quadratic companion error. In the quotient
just taken, brackets involving even lifts vanish, so an ordered
companion product becomes its additive characteristic-polynomial
expression. A derived element has the same degree-two logarithmic
and Magnus coefficient.

Finite generation clears the denominators of the chosen isomorphism
$U[1/p]\simeq R^r$ and its inverse: for some $a$,
$p^a\Lambda^r\subset U\subset p^{-a}\Lambda^r$ in these
coordinates. Completed tensor and exterior squares inherit fixed
denominator bounds. Consequently $N[1/p]$ is the bounded-denominator
alternating two-variable coefficient module, rather than a module of
series with arbitrary coefficient denominators.

Pass to a finite splitting field $L/\mathbb Q_p$, with integers
$\mathcal O$, and suppose that every coefficient of $C_e$ vanishes
on $D_\alpha=0$. All roots satisfy $\alpha-1\in\mathfrak m_{\mathcal O}$,
because the action of the pro-$p$ group $\Gamma$ on the finite-rank
lattice $T$ is continuous. The change of variable
\[
 S=(1+T_1)(1+T_2)-\alpha,\qquad
 T_2=\frac{\alpha+S}{1+T_1}-1
\]
is an integral formal-coordinate change. Its specialization kernel
is $(S)$, and division preserves a common denominator bound.
Since $D_\alpha$ is symmetric, an alternating numerator has an
alternating quotient. Thus, in $N_L$,
\[
 C_e=(\gamma-\alpha)A.                                      \tag{CE.4}
\]
Wrong-character coefficients are identically zero, so it suffices
to assume vanishing for the matching odd-character pairs.

Set $r_e(z)=q_e(z)/(z-\alpha)$. Cyclicity implies that
$w=r_e(\gamma)v$ is nonzero and spans an $\alpha$-eigenline
$W\subset T_L$. For a repeated root it is the socle of the cyclic
generalized eigenspace. Equations (CE.3)--(CE.4) show that
\[
 \widetilde w=r_e(\gamma)y-A
\]
is a $\Gamma$-equivariant lift of $w$. Hence (CE.2) splits over
$W$ after inverting $p$.

We now compare this split with the actual cup map. The lattice
$T_{\mathcal O}\cap W$ is saturated, hence an $\mathcal O$-direct
summand of $T_{\mathcal O}$. Choose an $\mathcal O$-linear section
$s:T_{\mathcal O}\to E_{\mathcal O}[1/p]$ extending the equivariant
section on that lattice. A fixed power $p^b$ makes this section
integral. Write the class-two logarithm of $h\in H$ as
\[
 x(h)+s(t(h))+z(h),\qquad
 x(h)\in U_{\mathcal O},\quad t(h)\in T_{\mathcal O},\quad
 z(h)\in N_{\mathcal O}[1/p].
\]
The same $p^b$ clears every $z(h)$. The section defect is
\[
 \beta_\delta(t)=\delta s(t)-s(\delta t),\qquad\delta\in\Gamma,
                                                               \tag{CE.5}
\]
and it vanishes on $W$.

Fix $n$ and two even cohomology classes $f,g$. Average their cocycles
by $(1+c)/2$. The resulting invariant cocycles vanish on the chosen
$\Gamma_n=\Gamma^{p^n}$ section: $c$ fixes its elements and is
$-1$ on the coefficient module. On $H$ they are additive maps
from $U$ to $\mathcal O/p^n(1)$, vanish on $T$, and are
$\Gamma_n$-equivariant. Put
\[
 \psi=f\wedge g:N_{\mathcal O}\longrightarrow\mathcal O/p^n(2),
 \quad
 \psi([x_1,x_2])=f(x_1)g(x_2)-f(x_2)g(x_1).
\]
The continuous cochain
\[
 k(h)=-\frac{p^b}{2}f(h)g(h)-\psi(p^b z(h))                    \tag{CE.6}
\]
is integral. The class-two product has central coordinate
$z_1+z_2+[x_1,x_2]/2$. Its substitution in (CE.6) gives
\[
 dk(h_1,h_2)=p^b f(h_1)g(h_2).
\]
Indeed the two summands contribute respectively
$p^b(f_1g_2+f_2g_1)/2$ and $p^b(f_1g_2-f_2g_1)/2$.

Extend $k$ to $P_n$ by $k(h\delta)=k(h)$. The cup has zero mixed
terms on the chosen section, so its Hochschild--Serre edge is, up
to the convention for the edge sign, the equivariance defect of
$k$. Its quadratic summand is equivariant; the remaining term is
\[
 \pm\psi(p^b\beta_\delta(t)).                               \tag{CE.7}
\]
This vanishes on $W$. A convention that conjugates $t$ makes no
difference because $W$ is $\Gamma$-stable. Thus the same fixed
$p^b$ multiple of every cup has zero restriction to the dual of
$W$, at every finite level.

For completeness, the target and its transitions are explicit.
Both $H$ and $\Gamma_n$ have $p$-cohomological dimension one,
and the minus part $U$ contributes only odd parity. Therefore
\[
 H^2(P_n,\mathcal O/p^n(2))^+
 \simeq H^1(\Gamma_n,\operatorname{Hom}(T,\mathcal O/p^n)(2)).
                                                               \tag{CE.8}
\]
Let $V=\operatorname{Hom}_{\mathbb Z_p}(T,\mathcal O)(2)$, with
$\gamma$ operator $B$. Normalized cocycle evaluation identifies
the right side with
\[
 V\big/\bigl(p^nV+(B^{p^n}-1)V\bigr).                         \tag{CE.9}
\]
Corestriction for the procyclic index-$p$ subgroup sends its
generator value to the same value in the larger group's
coinvariants. Thus the transitions are quotient maps, without an
extra factor of $p$. Continuity implies $B^{p^n}-1\to0$; the
denominators in (CE.9) are cofinal with the $p$-adic neighbourhoods.
Their inverse limit is $V$.

Restriction to the primitive lattice in $W$ gives, after inverting
$p$, a nonzero quotient
\[
 V_L\longrightarrow W^*(2),\qquad
 \gamma|_{W^*(2)}=u^2/\alpha,
\]
of character $\chi^2e^{-1}$, or exponent $2-e$ in cyclotomic
character notation. The restriction maps on the integral Hom
modules are surjective, and the corresponding finite coinvariant
quotients form a separated inverse system. By (CE.7) and the
finite-level cup-surjectivity hypothesis, the image of $p^bV$ in
every one of these quotients is zero. Its image in $W^*(2)$ must
therefore be zero. This contradicts the nonzero torsionfree
quotient. No compatible choice of finite-level cup expressions is
needed, since $b$ was fixed before their selection. The contradiction
proves (CE.1) for every root, including repeated ones.
\end{proof}
